% !PRINTED 122
\weberpage{122}
%<<CONT

nur endliche Werthe hat, und die ausserdem in diesem Intervall \gesperrt{stetig} ist; das will, in Uebereinstimmung mit \S.\,31, V und VI, besagen, dass
\begin{equation}
 f(x\pm h)-f(x)
\tag{2}
\end{equation}
dem absoluten Werthe nach in dem ganzen Intervall unter einer beliebig gegebenen Zahl $\eta$ liegt, wenn das positive $h$ kleiner als eine gewisse Zahl $\varepsilon$ ist. (Für $x=a$ nehmen wir in (2) nur das obere, für $x=b$ nur das untere Zeichen, um mit $x\pm h$ nicht aus dem Intervall herauszukommen.)

Für die Werthe einer solchen Function in dem Intervall giebt es nach 1. eine \gesperrt{untere Grenze} $g$, und wir beweisen nun den Satz:

3. \gesperrt{dass die Function $f(x)$ den Werth $g$ für irgend einen Werth $\xi$ des Intervalles annimmt, wonach die untere Grenze zu einem Minimum der Function $f(x)$ in dem Intervall wird.}

Der Beweis ist folgender. Die \gesperrt{untere Grenze} der Functionswerthe in einem Theil des Intervalles (1) ist nach 2. entweder gleich oder grösser als $g$.

Wenn nun $f(a)=g$ ist, so ist das zu Beweisende richtig; ist aber $f(a)>g$, so kann man wegen der Stetigkeit von $f(x)$ ein Intervall $a\le x\le a+h$ angeben, in dem $f(x)>g$ bleibt und also die \gesperrt{untere Grenze} von $f(x)$ grösser als $g$ ist.

Wir construiren nun in dem Intervall (1) einen Schnitt $(\mathfrak A,\mathfrak B)$ der Art, dass wir einen Werth $\alpha$ des Intervalles (1) zu $\mathfrak A$ rechnen, wenn die \gesperrt{untere Grenze} von $f(x)$ in dem Intervall $a\le x\le \alpha$ grösser als $g$ ist, und einen Werth $\beta$ zu $\mathfrak B$, wenn die \gesperrt{untere Grenze} im Intervall $a\le x\le \beta$ gleich $g$ ist.

$\mathfrak B$ kann möglicherweise aus dem einzigen Werthe $b$ bestehen; dann aber muss $f(b)=g$ sein, und unsere Behauptung ist für $x=b$ erfüllt. Denn wäre $f(b)>g$, so könnte man eine Grösse $g'$ zwischen $f(b)$ und $g$ und wegen der Stetigkeit ein Intervall $b-h\le x\le b$ so annehmen, dass in diesem Intervall alle Functionswerthe $f(x)$ grösser als $g'$, ihre \gesperrt{untere Grenze} also gleich oder grösser als $g'$ und daher sicher grösser als $g$ wäre. Da aber $b-h$ zu $\mathfrak A$ gehört, so ist auch in dem Intervall $a\le x\le b-h$ und folglich in dem ganzen Intervall (1) die \gesperrt{untere Grenze} grösser als $g$, gegen die Annahme.
